% ===============================================================
\chapter{Strawson on `Referring'}
% ===============================================================

\textbf{Strawson}: ``Clever guy'' (Eric Snyder). A believer in Ordinary Language Philosophy.

\section{Strawson's three-way distinction}

Strawson's three-way distinction that forms propositions:
\begin{enumerate}[label=(\alph*)]
  \item \emph{A literal sentence/expression.}
  \item \emph{The utterance of a sentence/\allowbreak expression}: the spatio-temporal context in which the sentence was uttered. It is not the utterance itself but the \emph{use} of the sentence (the statement made) that has a truth value; two different uses can make two different statements. When talking about utterances, he is talking about using sentence tokens, but when using ``the sentence'' he is talking about sentence types.
  \begin{enumerate}[label=\roman*.]
    \item Because the use has a truth value, the tokens of the sentences determine the truth value rather than the type.
    \item Expression types have meaning whether or not their presuppositions are met: meaning belongs to the type (the rules and conventions for its use). Presupposition failure only affects whether a use of the sentence is true or false (Strawson's criterion of assertion states that an assertion by a sentence is only true or false if all presuppositions are satisfied).
    \begin{itemize}
      \item This implies that Strawson does not use the principle of bivalence, which states that sentences can only be mapped to a semantic value of ``true'' or ``false''.
    \end{itemize}
    \item Meanings are given by the conventions and rules determined by how we live our everyday lives.
  \end{enumerate}
  \item \emph{The use of a sentence/\allowbreak expression}: the ``uniquely referring use'' of a particular sentence/\allowbreak expression to attach to a unique referent.
\end{enumerate}

\section{Neither true nor false statements}

If an expression like ``the present king of France'' is present in a proposition that was uttered now, the proposition is neither true nor false, because there is no referent for the present king of France, the key word being \emph{present}.

\section{Indexicals}

\begin{definition}
\emph{Indexicals} are linguistic expressions whose reference can shift from context to context. These are words like `I', `you', `here', `there'.
\end{definition}

Given how Strawson defines utterances and uses, two utterances of the same sentence type containing indexicals never have the same use. This is because they shift from context to context, and therefore have a different usage with a different time and spatial context.

\section{Entailment and presupposition}

\begin{definition}[Entailment]
A relation holding between (sets of) propositions, whereby $p$ entails $q$ if it is impossible for $p$ to be true while $q$ is false.
\begin{itemize}
  \item If $p \Rightarrow q$, then $p$ entails $q$.
\end{itemize}
\end{definition}

\begin{definition}[Presupposition]
A different relation holding between (sets of) propositions. Very roughly, $p$ presupposes $q$ if anyone sincerely asserting $p$ believes, and expects their interlocutors to also believe, $q$.
\begin{itemize}
  \item If $p \Rightarrow q$ and $\neg p \Rightarrow q$, then $p$ \emph{presupposes} $q$.
\end{itemize}
\end{definition}

There is also the family of sentences test:

\begin{slide}
\textbf{The Family of Sentences Test.} A sentence $S$ presupposes a sentence $S'$ if (i) $S$ takes for granted $S'$ and
\begin{enumerate}[label=(\roman*),start=2]
  \item `It's not the case that $S$' takes for granted $S'$
  \item `If $S$, then \ldots' takes for granted $S'$
  \item `It might be that $S$' takes for granted $S'$
  \item `Is it the case that $S$?' takes for granted $S'$
\end{enumerate}
\end{slide}

\section{Strawson's account of meaning}

Definite descriptions \emph{presuppose} existence. They do \emph{not entail} existence.

\begin{example}\leavevmode
\begin{enumerate}[label=(6\alph*)]
  \item The king of France is wise.
  \item There is a king of France.
\end{enumerate}
\end{example}

In light of this, why does Strawson write the following?

\begin{slide}
\itshape
There are at least two true things which Russell would say about the sentence [in (6a)]:
\begin{enumerate}[label=(\arabic*)]
  \item The first is that it is significant; that if anyone were now to utter it, he would be uttering a significant sentence.
  \item The second is that anyone now uttering the sentence would be making a true assertion only if there in fact at present existed one and only one king of France, and if he were wise.
\end{enumerate}
What are the false things which Russell would say about the sentence [in (6a)]? They are:
\begin{enumerate}[label=(\arabic*)]
  \item That anyone now uttering it would be making a true assertion or a false assertion.
  \item That part of what he would be asserting would be that there at present existed one and only one king of France.
\end{enumerate}
\end{slide}

\begin{definition}[Strawson's Criterion of Assertion]
An utterance of a sentence $S$ makes an assertion which is true or false only if all of $S$'s presuppositions are satisfied. Otherwise, an utterance of $S$ is neither true nor false.
\end{definition}

\section{Strawson's critiques of Russell}

\subsection*{Objection 1: The bearers of truth and falsity}

Russell maintains that:
\begin{premises}
  \item[P1)] Every grammatical sentence type expresses a \emph{unique} proposition.
  \item[P2)] Every proposition is either true or false.
\end{premises}

Strawson responded with:
\begin{enumerate}
  \item Without additional context/presuppositions, we cannot know if a sentence is true or false, which contradicts P2).
  \item Consider Strawson's example (7): `This is a fine red one.' In different contexts, different assertions are made, and those assertions may be true or false, depending on what the speaker is doing and whether it is red. But then the sentence type (7) cannot be said to express something true or false, but rather an utterance of that sentence, or rather a speaker's use of a token of that sentence type to make an assertion. This contradicts P1).
\end{enumerate}

\subsection*{Objection 2: Incomplete descriptions}

Why Russell's theory of descriptions collapses for definite descriptions. Consider

\begin{example}
(8) The table is covered with books.
\end{example}

\begin{argument}\leavevmode
\begin{premises}
  \item[P1)] If Russell's theory of descriptions is correct, then (8) would be true only if there is exactly one table in the universe.
  \item[P2)] There is more than one table in the universe.
  \item[C1)] So, if Russell's theory of descriptions is correct, then (8) must be false.
  \item[P3)] (8) needn't be false: in fact, it may be true.
  \item[C2)] Therefore, Russell's theory of descriptions is not correct.
\end{premises}
\end{argument}

There are two potential responses: the \emph{elliptical response} and \emph{quantifier domain restriction}.

\paragraph{Elliptical response.}
\begin{enumerate}
  \item Definite descriptions like `the table' are actually elliptical for more complex descriptions like `the table located in Room AC201 of Ashoka University \ldots'. In other words, `the table' in (8) really means `the table \ldots'. So long as there is enough elided material in `\ldots' to uniquely pick out a single table, (8) may be true.
  \item Think of the ellipses as being \emph{silent}.
  \item Problems with the elliptical response:
  \begin{enumerate}[label=(\alph*)]
    \item The elliptical material can be infinite, and can be further complex due to the definite description of the table containing multiple other definite descriptions.
    \item Sentence types with and without the definite description are distinct. This means that both sentences express the same proposition, which contradicts Russell's view that they must express unique ones.
  \end{enumerate}
\end{enumerate}

\paragraph{Quantifier domain restriction.}
\begin{enumerate}
  \item The quantifiers within a sentence essentially cover only a certain domain, which allows us to say ``the table'' after setting up a domain that is not the universe.
  \item Problems with this for Russell's theory:
  \begin{enumerate}[label=(\alph*)]
    \item Since we are doing all the work of restricting the domain, we don't really need Russell's theory.
    \item Russell wouldn't be comfortable with the truth/falsity of a sentence fluctuating across different domains.
    \item ``If a bishop meets another bishop, the bishop blesses the other bishop.'' Good luck restricting the domain for this.
  \end{enumerate}
\end{enumerate}
